\section{Discussion}
The results provided through the experimentation described in the previous section allow us to determine the effect of interaction style on the overall level of enjoyment, or flow, experienced while playing a game. In addition, the investigation of genre and the subject characteristics of gaming and touch-device experience allow for a greater understanding of the contributing factors related to this experience. This section will discuss and contrast the results obtained with regard to determining whether there is a difference in the level of flow experienced when using a gestural, direct manipulation or mixed interface. In addition, the effect of game genre and subject characteristics on this factor will be reviewed. Finally, a discussion of the results specific to the Puzzle game genre will be included.

\subsection{Effects of game genre on interface choice}
The significant interaction between interface and game genre suggests that the genre does effect the degree to which flow is elicited by each interaction style. However, the nature of this effect is less clear. All of the data showing a significant interaction between interface type and genre either referred to some effect between a genre-interface combination and another interface type within the same genre or a combination that included a different interface and genre type. In other words, no significant interactions were found between genre across the same interface type. In addition, when a significant difference between two interface types within the same genre was found, either other genres showed the same differences or no other genre showed a significant difference between any of their interface types. This is illustrated in the \emph{Challenge-Skill Balance} dimension, where within both the Fighting and Puzzle games, significant differences were found between direct manipulation and the other two interfaces. These differences were in the same direction, however. All other flow factors which show significant differences between interface types, exhibit these differences within one genre only, while the other genres show no significant between-interface differences. Furthermore, when a significant difference was found between two interfaces within a genre, both the interfaces identified and the direction of the interaction were corroborated by other significant findings within the \emph{Interface} main effect as well as the \emph{Interface}*\emph{Group} interaction.

However, a noticeable discrepancy is present when considering the results of this interaction for the higher-order factor. That is, the direct manipulation interface appears to elicit more flow than either of the other two versions for the Strategy game. This contradicts the findings of a number of individual flow factors. While the difference between the direct manipulation and gestural interface is not significant, direct manipulation appears to elicit significantly more flow than the mixed interface along this higher-order factor. As this factor is an aggregate of the nine individual flow factors, this may be due to the significant differences found between these interface version within the \emph{Loss of Self-Consciousness}, \emph{Transformation of Time} and \emph{Autotelic Experience} dimensions. These significant differences appear to have greatly affected the results of the Strategy game with regard to the higher-order factor. However, as this factor is a combination of the results of the individual flow factors, it may be more appropriate to examine their results, rather than their aggregate, in order to obtain a more fine-grained understanding of how genre affects interface preference.

The ordinal interaction present across most flow dimensions, also suggests that the effect of genre is consistently in the same direction. In other words, the interface versions elicit similar levels of flow across game genre. In addition, the almost parallel nature of these lines would suggest that this interaction is fairly weak. This ordinal interaction is especially true when considering the Strategy and Fighting/Duelling games. The nature of the Puzzle game may explain why it does not as closely match the means of the previous two genres. That is, it may be more difficult to elicit flow, or the levels along which flow was elicited may have varied, with such a game, when compared to Strategy or Fighting/Duelling games. One of the characteristics present in the other two genres but not in the Puzzle game, were real-time constraints placed on the user during game-play. In the Puzzle game, unlike in the other two game genres, where issues of timing had to be considered, the user was able to spend as long as they liked on any particular element. This may have led to the participant either not experiencing flow as strongly or simply eliciting flow along different dimensions to those of the other two games. Similar reasoning may be used to explain the disordinal interaction present for the mixed condition when considering the Strategy game. Furthermore, when removing the mixed condition, there is agreement in the distribution of means between all three genres. That is, the interactions between direct manipulation and the gestural interface are ordinal across the three genres for all flow factors when considering the results of the MANOVA. The discrepancies introduced by the mixed interface version may be due to the differences in the ways in which direct manipulation and gestural commands were combined within each genre. That is, while in the other two interface versions all commands were either gestural or direct manipulation, the manner in which these two interaction techniques were combined in the mixed version may have resulted in additional variance related to their combination rather than the actual genre.

It should also be noted that any discrepancies may be specific to the actual games, rather than the genre as a whole. As there is considerable variation in game design, even within a game genre, additional testing would be required to definitively say that these differences are due to game genre as a whole, rather than just the specific attributes of these individual games. Thus, it appears that, while genre does affect the degree to which flow is experienced, this effect is not easily explained. Genre appears to accentuate or dampen the degree to which flow is experienced within interface types along specific flow dimensions. However, these variations are not large or consistent enough to suggest that genre plays a definitive role in which interface type is preferred. Rather, the presence of a significant interaction seems to suggest that genre effects the degree to which flow is experienced within the larger overall direction shared by all genres, rather than serving as a primary contributing factor in the preference of interaction technique.

\subsection{Interaction between participant characteristics and interaction style}
The effect of previous gaming and touch-device experience was explored in relation to interaction style preferences. A significant interaction was found between these two variables. However, as with the \emph{Interface}*\emph{Genre} interaction, understanding the nature of this interaction is made more complex by the number of dependent variables that were assessed. As with the \emph{Interface}*\emph{Genre} interaction, no significant differences were found within an interface type across the participant characteristic groups. However, differences were found within specific groups between different interface versions, but this was only present in the Gaming and Touch-device experience conditions. In other words, the Control group showed no significant differences between interface types.

Where significant differences were found, both the nature and direction of these differences were consistent with the other significant interactions found for that flow dimension. In addition, when differences were present between interface types within both the Gaming and Touch-device experience groups for a specific flow factor, the direction and interface types affected were shared. However, significant results for the Touch-device experience group were only obtained for flow dimensions where the same was found for the Gaming experience condition. The converse was not always the case: the gaming experience group showed far more significant between-interface differences than either of the two other groups. As a significant result implies that the means of the two conditions were sufficiently far apart, this may mean that those in the Gaming experience condition experienced or did not experience flow more strongly than those in the other two experimental groups. That is to say, when a flow dimension elicited large amounts of flow, those in the Gaming experience group experienced this to a greater degree than those in the other two groups, while the same is true of the converse. That is, when flow was only weakly experienced, those in the Gaming experience group experienced less flow than individuals with only limited gaming experience. Essentially, gamers seem to experience an amplification of flow both positively and negatively. This may be due to a number of factors. Firstly, due to previous gaming experience, these individuals may be able to more readily enter a flow state than those in the other two conditions. In addition, they may thus be more sensitive to elements that disrupt this flow state. A second explanation relates to personality factors that predispose certain individuals to being more amenable to entering a flow state~\citep{Csik90}. That is, those with gaming experience may have a personality type that predisposes them to entering a flow state more readily, and because of this, enjoy playing games more so than those who lack this trait. This in turn may explain why they play games more often than other individuals and also why they are more sensitive to variations in flow when compared to those in the Touch-device experience or Control groups. However, further testing would be required to definitively accept either hypothesis as the cause of this variation.

The Multivariate contrast analysis showed significant differences between the Control and other groups. The mean comparison for this test seems to suggest that possessing either Gaming or Touch-device experience does increase the degree of flow experienced overall. That is, individuals who possess these participant characteristics more readily enter a flow state than those who have both limited gaming and touch-device experience, regardless of interface type. However, as those in the Touch-device experience group had more gaming experience (as assessed by the pre-experiment questionnaire) than those in the Control group, gaming experience may still be the active factor in this case. Without further experimentation, however, it is not possible to further separate these variables. With this said, the results are not surprising overall, as an unfamiliarity with either an interaction technique or an application-type may cause frustration as well as feelings of confusion, which in turn, would decrease the amount of flow experienced, regardless of interface type. 

Thus, while interface preference appears to have been fairly consistent both across groups, as well as when compared to the significant results for the \emph{Interface} main effect and \emph{Interface}*\emph{Genre} interaction, participant characteristics do appear to have influenced the flow results. It seems that those with touch-device or gaming experience are more likely to enter a flow state of a higher degree across most flow dimensions, regardless of interface type. In addition, individuals who regularly play games appear to be more sensitive to overall variations present for each interface type within the flow dimensions. Thus, when considering interface type, participant characteristics do appear to influence the degree of the flow state experienced, rather than its specific nature.

\subsection{Flow and the effect of interaction technique}
Due to the nature of the data, it is fairly difficult to make clear statements with regard to which interaction technique elicits more flow. Despite this, a number of patterns seem to have emerged. As the nature of the interactions between \emph{Interface} and both \emph{Genre} and \emph{Group} have been reviewed above, only the results specific to differences in flow between the various interfaces (including when these differences contain interactions with other variables) will be discussed in this section. 

A number of dimensions indicated a clear preference for the gestural and mixed interface over the direct manipulation version. These include \emph{Challenge-Skill Balance}, \emph{Merging of Action and Awareness} and \emph{Clear Goals}. The preference in the \emph{Challenge-Skill Balance} dimension may indicate that the addition of gestural controls affects the level of difficulty present in playing the games. However, without further testing it is not possible to specify whether this increases or decreases the difficulty of the game. The results for \emph{Merging of Action and Awareness} seem to suggest that interfaces containing gestural elements increase the degree to which immersion is experienced. This is to be expected, in a sense, as the combination of a touch-based interaction technique with more fluid gestural controls may remove some of the perceived distance present between the user, device and application. The explanation for the preference of the gestural and mixed interfaces in the case of \emph{Clear goals} is less apparent, however. As all interface versions shared the same goals within each game genre, this difference seems to suggest that these goals were more apparent when using gestural or a combination of gestures and direct manipulation. Finally, the mixed version does not elicit a significantly greater flow state along any other flow dimensions. The presence of differences only in cases where the same was found between the gestural and direct manipulation interfaces suggests that the source of the higher flow state in the mixed version may be the inclusion of gestural commands. 

The gestural interface was shown to elicit more flow than the direct manipulation interface along the \emph{Unambiguous Feedback}, \emph{Concentration on the Task at Hand} and \emph{Sense of Control} dimensions. This result, in the case of \emph{Unambiguous Feedback} and \emph{Sense of Control}, is unexpected. Due to the general unfamiliarity of users with regard to gestures, one might expect that the direct manipulation interface would provide users with more of a feeling of control. The mixed interface was also shown to elicit significantly less flow than the gestural version along this dimension. In this case, one interpretation is that the use of direct manipulation controls within this interface may have negated any positive effects associated with the included gestural controls. A second explanation is that the inclusion of direct manipulation controls as well as gestural ones may have confused users. With regard to \emph{Unambiguous Feedback}, as the feedback given when a command was executed was the same across interface type, this difference suggests that the use of gestures enhanced the degree to which users were able to interpret this information. However, the mechanism behind this is unclear and would require further testing. Finally, as gestures have been shown to increase the cognitive load placed upon users~\citep{Niel04}, they may have had to concentrate more in the gestural version than in the other two in order to successfully interact with the game. This would potentially explain the higher flow state experienced in the gestural version when considering \emph{Concentration on the Task at Hand}.

The \emph{Loss of Self-Consciousness} dimension showed a significantly lower flow score for the mixed interface than either of the other two versions. No significant interaction was found between the direct manipulation and gestural interfaces, however. This suggests that the combination of interaction techniques plays some role in decreasing the degree to which users are able to disassociate from themselves and immerse themselves in the game-play. Once again, the reasons for this are unclear and a deeper understanding of this phenomenon would require further investigation. A similar effect was found for \emph{Transformation of Time}, except in this case only the direct manipulation interface was shown to elicit significantly more flow than the mixed version. Again, the potential confusion and possible additional cognitive load present when combining multiple interaction techniques should perhaps be considered.

Finally, with regard to \emph{Autotelic Experience}, direct manipulation elicited significantly more flow than either the gestural or mixed interfaces. This is surprising, as one might expect that the use of gestures would provide more of an intrinsically rewarding activity due to the process involved in executing each command. However, the increased cognitive load as well as learning requirements present in the gestural case may have influenced this. That is, in the direct manipulation version, the user is able to play the game with only minimal game-play related learning required. In the gestural interface, the amount of learning required is compounded by the use of gestures. This, in turn, may detract from the intrinsic enjoyment value of the game. Increased familiarity with the gestural commands might lessen this difference if it is indeed due to the cognitive demands placed on the user by the gestural interface.

Overall, further experimentation is required when attempting to explain the differences between the various interfaces on a dimension by dimension basis. While possible explanations for these discrepancies have been presented above, they are not conclusive. In addition, these factor results do not provide an overall view of the appropriateness of interaction techniques within gaming. They do, however, present an interesting pattern.

That is, it appears that the gestural interface elicits flow more highly and consistently than either of the other two interfaces. This is further confirmed by the Multivariate contrast analysis. This analysis also showed that the mixed interface version was preferred to the direct manipulation one. In addition, the mixed interface does not elicit significantly more flow than the direct manipulation interface in any conditions except where the same is true of the gestural interface. These two factors may suggest this effect is due to the inclusion of gestures in the mixed interface. The fact that direct manipulation was felt to be more intrinsically rewarding seems to contradict the other results, however. This suggests the need for further testing to determine whether this is due to an unfamiliarity with gestural techniques, and an increase in the cognitive load placed on the user. Overall, however, when considering gaming experience, it appears that gestures are the most successful at eliciting a higher flow state.

The consistency of results regardless of genre or participant characteristics suggests that it was indeed interface type which determined the way in which flow was experienced in these games. Within this, the inclusion of gestures in an interface appears to increase the level of flow experienced by participants. Due to the dimensionality present in the results, these findings should be interpreted cautiously. However, the consistency with which the gestural interface elicits more flow, as well as the results of the contrast analysis, suggest that gestures do indeed provide the most enjoyable gaming experience in this case.

\subsection{Results specific to the Puzzle game}
The Multivariate contrast analysis conducted on the results of the Puzzle game provided information that corresponds with the data described above. This is unsurprising given the lack of any large inter-genre differences when examining interface preferences.

Within this genre, the gestural interface was preferred over either of the other two interaction techniques, although no significant difference was found between the direct manipulation and mixed interfaces. Given that the game revolved around a three-dimensional shape, the preference for gestural commands may have been due to the fact that they allowed the user to feel as if they were interacting with the shape on a more physical level. This feeling would, in turn, have increased the degree of immersion experienced. As a level of immersion is required for flow to be experienced~\citep{Swee05}, this would have increased the overall flow state experienced. 

The contrast analysis also showed a significant difference between the Gaming and Control groups. As discussed above, this may be due to either the personality characteristics of those who play games frequently, or due to their more extensive experience in entering a flow state.

Perhaps the most important feature of these results, however, is that they correlate with the results obtained overall. Given that the Puzzle game exhibited a number of distinct design differences when compared to the other two genres, this reaffirms the hypothesis that genre does not overtly influence interface preference. These differences include the exclusion of elements such as real-time constraints as well as a difference in the actual aim of the game when compared to the other two genres. That is, within the Puzzle game, the player attempted to achieve as high a score as possible, while the other two genres required the player to defeat some opponent. 

The fact that the results of the Puzzle game corroborate those achieved for all three games suggests that genre, and more specifically, differences in game-play and game aims, do not play an important role in the degree to which flow is experienced with regard to interaction style. Moreover, the lack of significant differences along this dimension further reaffirms the idea that the flow differences present are indeed due to the choice of interaction technique.



